\chapter{VVC}
\index{VVC}

\section*{Definition and Overview}
Vulvovaginal candidiasis (VVC), colloquially referred to as a vaginal yeast infection or vaginal thrush, is a highly prevalent mucosal infection of the female lower reproductive tract. It is characterized by acute inflammation of the vulva and vaginal mucosa, resulting from an overgrowth of opportunistic fungal species belonging to the genus \textit{Candida}. Under physiological conditions, the female lower genital tract maintains a delicate immunological and microbiological equilibrium. Fungi of the genus \textit{Candida} reside as benign commensals in the vaginal lumen of approximately 20\% to 30\% of healthy, asymptomatic individuals. However, when host physiological defenses are compromised, or when the vaginal microenvironment is disrupted, these organisms can transition from a harmless colonizing state to an active pathogenic state, leading to symptomatic disease.

While VVC is not classified as a sexually transmitted infection (STI) because the causative organisms are endogenous and the condition occurs in sexually inactive individuals, its clinical presentation frequently overlaps with common STIs. This clinical overlap necessitates a rigorous diagnostic approach to prevent misdiagnosis and inappropriate treatment. In clinical settings, VVC is categorized into two main classifications: uncomplicated VVC and complicated VVC. Uncomplicated VVC is defined as sporadic, infrequent episodes of mild-to-moderate severity occurring in immunocompetent individuals, typically caused by the highly susceptible species \textit{Candida albicans}. Complicated VVC refers to episodes that are severe, recurrent (defined as four or more symptomatic episodes within a twelve-month period), caused by non-\textit{albicans} species, or occurring in individuals with predisposing systemic conditions such as poorly controlled diabetes mellitus or immunosuppression. Understanding this distinction is vital, as it guides the clinical management, drug selection, and treatment duration required to achieve therapeutic success and minimize patient morbidity.

\section*{Epidemiology}
The epidemiological burden of vulvovaginal candidiasis is substantial, representing one of the most common reasons for gynecological consultations, primary care visits, and self-medication globally. It is estimated that approximately 70\% to 75\% of individuals with a uterus will experience at least one episode of symptomatic VVC during their lifetime, typically during their reproductive years. Within this cohort, approximately 40\% to 50\% will suffer from a second or subsequent episode, highlighting a high rate of recurrence. Furthermore, a smaller but clinically challenging subset of the population, estimated at 5\% to 9\% of adult women, experiences recurrent vulvovaginal candidiasis (RVVC). RVVC represents a chronic and debilitating condition that imposes a severe physical, emotional, and financial burden on affected individuals, often requiring long-term suppressive therapy and lifestyle modifications.

VVC is closely linked to reproductive physiology and ovarian hormone levels, which explains its distribution across different demographics. The incidence of VVC rises sharply after menarche, peaks during the second and third decades of life, and declines significantly following menopause, unless the individual is receiving postmenopausal hormone replacement therapy (HRT) or has an underlying metabolic disorder such as diabetes. Prepubertal girls are rarely affected by VVC; the presence of symptoms in this age group should prompt investigations for other causes, such as foreign bodies, poor hygiene, or potential abuse. While the condition occurs across all socioeconomic and ethnic groups, certain behavioral and environmental factors can influence prevalence. These include the use of specific contraceptive methods (e.g., high-dose combined oral contraceptives or intrauterine devices), sexual activity patterns, and access to healthcare resources for diagnosis and management.

\section*{Aetiology and Pathophysiology}
The development of symptomatic vulvovaginal candidiasis involves a complex interaction between the fungal pathogen, the vaginal microbiome, host mucosal immunity, and predisposing systemic factors. The transition of \textit{Candida} from a harmless commensal to an invasive pathogen is driven by changes in host physiology and fungal virulence factors.

Key factors in the etiology and pathophysiology of VVC include:
\begin{itemize}
    \item \textbf{Causative Organisms}: The primary pathogen in VVC is \textit{Candida albicans}, accounting for approximately 85\% to 90\% of all cases. The remaining 10\% to 15\% of cases are caused by non-\textit{albicans} \textit{Candida} (NAC) species. Among these, \textit{Candida glabrata} (sometimes classified as \textit{Nakaseomyces glabrata}) is the most common, followed by \textit{Candida tropicalis}, \textit{Candida parapsilosis}, \textit{Candida krusei}, and \textit{Candida lusitaniae}. Non-\textit{albicans} species are of particular clinical significance because they are often less susceptible or intrinsically resistant to standard azole antifungal therapies, frequently leading to chronic, recurrent, or treatment-refractory cases.
    \item \textbf{Morphological Transition}: \textit{Candida albicans} is a dimorphic fungus capable of transitioning between two primary morphological forms: the unicellular budding yeast (blastospore) form and the multicellular filamentous (hyphal or pseudohyphal) form. While the yeast form is associated with colonization and dissemination, the hyphal form is highly invasive. Hyphae produce specialized cell-surface proteins (adhesins) that bind to vaginal epithelial cells and secrete hydrolytic enzymes, such as Secreted Aspartyl Proteinases (SAPs) and phospholipases, which degrade host cellular membranes and facilitate tissue penetration.
    \item \textbf{Candidalysin and Epithelial Damage}: During the hyphal transition, \textit{Candida albicans} secretes a cytolytic peptide toxin called candidalysin, encoded by the \textit{ECE1} gene. Candidalysin damages the host cell membrane by forming pores, leading to the leakage of intracellular contents and triggering an inflammatory signaling cascade in the vaginal epithelial cells.
    \item \textbf{Host Immunopathogenesis}: The symptoms of VVC are not primarily caused by the direct destruction of tissue by the fungus, but rather by an intense, hyper-inflammatory host immune response. The epithelial cell damage caused by candidalysin induces the release of pro-inflammatory cytokines, particularly interleukin-1 beta (IL-1 beta), and chemotactic chemokines. This chemical signaling recruits a massive influx of polymorphonuclear neutrophils (PMNs) to the vaginal mucosa. These neutrophils, however, fail to clear the fungal burden effectively and instead release reactive oxygen species and lysosomal enzymes, driving the local tissue inflammation, pruritus, and pain characteristic of VVC.
    \item \textbf{Vaginal Microbiome and Lactobacilli}: A healthy vaginal ecosystem is dominated by \textit{Lactobacillus} species (such as \textit{L. crispatus}, \textit{L. gasseri}, \textit{L. jensenii}, and \textit{L. iners}). Lactobacilli maintain a low vaginal pH (typically between 3.8 and 4.5) through the production of lactic acid. This acidic environment inhibits the yeast-to-hyphae transition. Lactobacilli also compete with \textit{Candida} for nutrients and adherence sites, and they produce hydrogen peroxide and bacteriocins that suppress fungal growth. Any disruption of this bacterial population can lead to \textit{Candida} overgrowth.
    \item \textbf{Broad-Spectrum Antibiotic Use}: The administration of systemic broad-spectrum antibiotics (such as tetracyclines, ampicillin, and cephalosporins) is a common trigger for VVC. Antibiotics eradicate the protective vaginal lactobacilli, removing the biological competition and allowing \textit{Candida} species to proliferate and cause symptomatic infection.
    \item \textbf{Hormonal Influences}: High estrogen states (such as during pregnancy, the luteal phase of the menstrual cycle, or the use of combined oral contraceptives) increase the glycogen content of vaginal epithelial cells. Glycogen serves as an abundant carbon source for fungal growth and acid production. Estrogen also increases the expression of vaginal cell receptors that bind \textit{Candida} adhesins.
    \item \textbf{Metabolic and Immunological Factors}: Poorly controlled diabetes mellitus leads to elevated glucose levels in vaginal secretions, which promotes fungal adherence and virulence. Immunodeficiencies, whether caused by HIV infection, chemotherapy, or immunosuppressive therapy, impair local cell-mediated immunity, making the host highly susceptible to severe and recurrent infections.
\end{itemize}

\section*{Clinical Features}
The clinical presentation of vulvovaginal candidiasis varies in severity, ranging from mild, localized irritation to severe, debilitating inflammation of the external genitalia and vaginal canal. Symptoms are typically localized and are often reported to worsen in the days preceding menstruation due to cyclical hormonal changes.

Common clinical features of VVC include:
\begin{itemize}
    \item \textbf{Vulvar Pruritus}: This is the most common and characteristic symptom of VVC, reported by almost all symptomatic individuals. The itching is often intense and persistent, leading to significant distress and sleep disruption.
    \textbf{Vaginal Discharge}: The classic discharge associated with VVC is thick, white, and clumpy, resembling cottage cheese or curdled milk. However, the discharge can also be thin, watery, or virtually absent in some cases. It is characteristically odorless, which helps distinguish it from other infections.
    \item \textbf{Soreness and Burning}: Patients frequently report a raw, sore, or burning sensation in the vulva and vagina, which is often exacerbated by urination (causing external dysuria) and by sexual intercourse (dyspareunia).
    \item \textbf{Erythema and Edema}: Physical examination typically reveals marked redness (erythema) and swelling (edema) of the vulva, labia minora, labia majora, and the vaginal mucosa.
    \item \textbf{Excoriation and Fissures}: The intense itching often leads to scratching, resulting in linear scratches (excoriations), superficial skin breakdown, and painful cracks in the skin (fissures), particularly in the interlabial folds and perineum.
    \item \textbf{Vaginal Plaques}: Visual inspection of the vaginal walls using a speculum often reveals adherent, white, patch-like plaques. When these plaques are wiped away, they reveal an inflamed, red mucosal base.
    \item \textbf{Atypical Presentations}: Infections caused by non-\textit{albicans} species, particularly \textit{Candida glabrata}, may present with atypical or milder symptoms, such as chronic vaginal burning or irritation without the characteristic thick discharge, which can lead to delayed diagnosis.
\end{itemize}

\section*{Differential Diagnosis}
Given that the symptoms of vulvovaginal candidiasis are non-specific and overlap with many other lower reproductive tract disorders, an accurate diagnosis cannot be made based on clinical history and physical examination alone. It is critical to differentiate VVC from other infectious and non-infectious conditions to ensure appropriate management.

Key differential diagnoses to consider include:
\begin{itemize}
    \item \textbf{Bacterial Vaginosis (BV)}: A polymicrobial syndrome characterized by a shift in vaginal flora from lactobacilli to anaerobic bacteria (e.g., \textit{Gardnerella vaginalis}). Unlike VVC, BV typically presents with a thin, homogeneous, greyish-white discharge that has a fishy odor, particularly after intercourse or the addition of potassium hydroxide (positive whiff test). Vaginal inflammation, itching, and pain are rare, the vaginal pH is elevated (greater than 4.5), and microscopy reveals clue cells and a lack of yeast.
    \item \textbf{Trichomoniasis}: A sexually transmitted infection caused by the protozoan parasite \textit{Trichomonas vaginalis}. It presents with a profuse, frothy, yellow-green vaginal discharge, vulvar irritation, dysuria, and malodor. Physical examination may reveal a ``strawberry cervix'' (punctate mucosal hemorrhages). The vaginal pH is elevated (greater than 4.5), and saline wet mount microscopy demonstrates motile, flagellated trichomonads and abundant neutrophils.
    \item \textbf{Contact Dermatitis (Irritant or Allergic)}: A non-infectious inflammatory skin condition of the vulva caused by exposure to local irritants (e.g., soaps, bubble baths, douches, laundry detergents) or allergens (e.g., latex, perfumes, topical medications). Symptoms include burning, itching, and redness, but vaginal discharge is absent, vaginal pH is normal, and microscopic examination shows no fungal elements.
    \item \textbf{Lichen Sclerosus}: A chronic, inflammatory skin disease that primarily affects the vulvar and perianal regions. It is characterized by intense vulvar pruritus, but unlike VVC, it leads to structural changes, including thinning of the skin (creating a white, shiny, ``parchment-like'' appearance), loss of the labia minora, and narrowing of the vaginal opening. It does not affect the vaginal mucosa. Diagnosis is confirmed via clinical evaluation and skin biopsy.
    \item \textbf{Genital Herpes Simplex Virus (HSV) Infection}: An STI caused by HSV-1 or HSV-2. Primary episodes present with painful, grouped vesicles on an erythematous base that rupture to form shallow, extremely tender ulcers. These local findings are often accompanied by systemic symptoms such as fever, myalgia, headache, and bilateral inguinal lymphadenopathy.
    \item \textbf{Physiological Leucorrhoea}: The normal, healthy vaginal secretion that varies in volume and consistency throughout the menstrual cycle. It is typically clear or white, non-irritating, and odorless. Microscopic examination shows normal epithelial cells, abundant lactobacilli, a pH of less than 4.5, and no fungal elements or pathogens.
    \item \textbf{Cytolytic Vaginosis}: A condition characterized by an overgrowth of lactobacilli, leading to excessive lactic acid production and cytolysis of vaginal epithelial cells. The clinical features, including pruritus, burning, and a white discharge, are virtually identical to VVC. However, microscopic examination reveals abundant lactobacilli, fragmented epithelial cell nuclei (naked nuclei), an extremely low vaginal pH (3.5 to 4.0), and a complete absence of \textit{Candida} species.
\end{itemize}

\section*{When to Seek Medical Care}
While vulvovaginal candidiasis is a common and generally non-life-threatening condition, obtaining an accurate professional diagnosis is essential to ensure appropriate treatment and rule out more serious pelvic or systemic infections. Individuals should consult a qualified healthcare provider, such as a general practitioner, gynecologist, or sexual health clinician, under the following circumstances:
\begin{itemize}
    \item This is the individual's first experience with vulvovaginal symptoms, as self-diagnosis without clinical validation is frequently incorrect and can delay the management of other conditions.
    \item Symptoms do not improve or resolve within three to five days of completing an over-the-counter antifungal treatment.
    \item Symptoms are severe, including intense swelling, severe pain, extensive skin splitting (fissures), or bleeding.
    \item The individual experiences frequent recurrences, defined as four or more infections per year.
    \item The individual is pregnant, as certain medications are contraindicated during pregnancy and special care must be taken.
    \item The individual has an underlying medical condition that compromises the immune system, such as poorly controlled diabetes, HIV/AIDS, or is undergoing chemotherapy.
    \item Symptoms are accompanied by systemic signs such as fever, chills, nausea, vomiting, lower abdominal or pelvic pain, or a foul-smelling vaginal discharge, which suggest a pelvic inflammatory disease or a systemic infection rather than localized VVC.
\end{itemize}

In the event of severe, life-threatening symptoms, such as high fever, severe pelvic pain, vomiting, or signs of systemic sepsis, contact emergency services immediately (e.g., local emergency services in many countries, 911 in the United States, or 999 in the United Kingdom).

\section*{Diagnosis and Investigations}
An accurate diagnosis of vulvovaginal candidiasis requires a combination of clinical assessment and laboratory investigations. Self-diagnosis and diagnosis based solely on clinical signs are unreliable, often leading to inappropriate use of antifungal agents and chronic vaginal symptoms.

Diagnostic investigations for VVC include:
\begin{itemize}
    \item \textbf{Clinical Examination}: A thorough pelvic examination, including visual inspection of the vulva and vagina and a speculum examination, is performed to assess the severity of inflammation, evaluate the characteristics of the discharge, and check for the presence of fissures, plaques, or signs of alternative diagnoses.
    \item \textbf{Vaginal pH Testing}: A sample of vaginal discharge is collected from the lateral walls of the vagina using a narrow-range pH paper. In uncomplicated VVC, the vaginal pH is characteristically normal (less than 4.5). An elevated pH (greater than 4.5) suggests bacterial vaginosis, trichomoniasis, or mixed infections, and helps narrow down the differential diagnoses.
    \item \textbf{Wet Mount Microscopy (KOH and Saline Prep)}: This remains the primary diagnostic tool in the clinical setting. A swab of vaginal discharge is mixed with a drop of normal saline on a glass slide, and another sample is mixed with a drop of 10\% potassium hydroxide (KOH). The KOH solution lyses host cellular debris (red blood cells, white blood cells, and epithelial cells), clearing the field to visualize the yeast cells, budding blastospores, and pseudohyphae or true hyphae. Saline preparation allows for the evaluation of clue cells and motile trichomonads. Microscopic visualization of fungal elements confirms the diagnosis in symptomatic individuals.
    \item \textbf{Fungal Culture}: Vaginal culture on Sabouraud dextrose agar or other specialized fungal media is not routinely required for patients with typical, uncomplicated VVC who respond to standard therapy. However, culture is indicated for:
    \begin{itemize}
        \item Patients with recurrent VVC (RVVC) to confirm the diagnosis and identify the species.
        \item Patients who fail to respond to standard antifungal therapy, to rule out non-\textit{albicans} species.
        \item Symptomatic patients who have normal vaginal pH and negative microscopic findings.
        \item Cases where susceptibility testing is needed due to suspected drug resistance.
    \end{itemize}
    \item \textbf{Molecular Testing (PCR)}: Polymerase chain reaction (PCR) assays and other nucleic acid amplification tests (NAATs) are increasingly available for the detection of \textit{Candida} DNA. These molecular tests have higher sensitivity and specificity than microscopy and culture and can identify multiple species simultaneously. However, because \textit{Candida} can be a normal commensal, a positive PCR test must always be interpreted in the context of the patient's clinical symptoms.
\end{itemize}

\section*{Management}
The management of vulvovaginal candidiasis is tailored to the severity of the clinical presentation, the frequency of episodes, the causative species, and the patient's physiological status (such as pregnancy or immunosuppression). The primary goal is the rapid alleviation of symptoms and the eradication of fungal overgrowth.

Treatment strategies include:
\begin{itemize}
    \item \textbf{Uncomplicated VVC Treatment}: Uncomplicated VVC responds well to short-course topical or oral antifungal regimens. Both routes of administration are equally effective, and the choice depends on patient preference.
    \begin{itemize}
        \item \textbf{Topical Azoles}: Creams, ointments, or vaginal suppositories containing clotrimazole (1\% or 2\%), miconazole (2\% or 4\%), or terconazole (0.4\% or 0.8\%) are administered intravaginally for 1 to 7 days. Topical regimens are preferred for patients who wish to avoid systemic medications or those with potential drug-drug interactions with oral agents.
        \item \textbf{Oral Antifungals}: A single oral dose of fluconazole (150 mg) is highly effective and widely prescribed due to its convenience. Therapeutic levels are maintained in vaginal secretions for up to 72 hours.
    \end{itemize}
    \item \textbf{Complicated VVC Treatment}: Complicated infections require more intensive and prolonged therapeutic courses.
    \begin{itemize}
        \item \textbf{Severe Inflammation}: Patients with severe vulvovaginitis (marked erythema, edema, and excoriation) should receive either two sequential doses of oral fluconazole (150 mg given 3 days apart) or a 7-to-14-day course of topical azole therapy. A mild topical corticosteroid (such as 1\% hydrocortisone cream) may be applied to the vulva for 2 to 3 days to alleviate severe pruritus and burning.
        \item \textbf{Pregnancy}: Only topical azole therapies (e.g., clotrimazole or miconazole cream or suppositories for 7 days) are recommended during pregnancy. Oral fluconazole is contraindicated, especially in the first trimester, due to a potential risk of congenital abnormalities and spontaneous abortion.
        \item \textbf{Non-albicans Infections}: Non-\textit{albicans} species, particularly \textit{Candida glabrata}, are relatively resistant to standard azole therapies. First-line treatment for symptomatic \textit{C. glabrata} VVC is intravaginal boric acid (600 mg in a gelatin capsule) administered daily at bedtime for 14 days. Other options include topical flucytosine or amphotericin B creams, which are typically compounded.
    \end{itemize}
    \item \textbf{Recurrent VVC (RVVC) Maintenance Therapy}: For patients suffering from RVVC, the goal is long-term suppression rather than a quick cure. The standard regimen consists of:
    \begin{itemize}
        \item \textbf{Induction Phase}: Oral fluconazole (150 mg) administered every 3 days for a total of 3 doses (days 1, 4, and 7) to achieve clinical remission.
        \item \textbf{Maintenance Phase}: Oral fluconazole (150 mg) taken once weekly for 6 months. Upon cessation of maintenance therapy, approximately 30\% to 50\% of patients will experience a recurrence, requiring re-evaluation or alternative regimens such as oteseconazole (a novel, highly selective oral CYP51 inhibitor approved for the prevention of RVVC in females who are not of reproductive potential).
    \end{itemize}
    \item \textbf{Non-Pharmacological and Supportive Care}: Patients should avoid local irritants and practice good vulvar hygiene. Partner treatment is not routinely recommended, as VVC is not considered an STI, and treating asymptomatic partners does not reduce recurrence rates in the index patient. Partner evaluation is only indicated if the partner is symptomatic (e.g., presenting with balanitis).
\end{itemize}

\section*{Complications}
Although vulvovaginal candidiasis is primarily a localized mucosal infection, it can lead to several physical and psychological complications, particularly if left untreated or in cases of chronic, recurrent disease.

Potential complications include:
\begin{itemize}
    \item \textbf{Secondary Bacterial Infection}: Intense pruritus leads to scratching, which damages the vulvar skin barrier. This can result in painful excoriations, superficial skin splits (fissures), and secondary bacterial infections (cellulitis or impetigo) caused by common skin pathogens like \textit{Staphylococcus aureus} or \textit{Streptococcus} species.
    \item \textbf{Chronic Vulvar Pain and Vulvodynia}: Repetitive, severe inflammatory episodes of VVC can sensitize local nociceptors, contributing to the development of chronic vulvar pain syndromes, such as provoked vestibulodynia (localized pain at the vaginal entrance upon touch or penetration).
    \item \textbf{Obstetric Complications}: In pregnant individuals, untreated VVC has been tentatively linked in some studies to an increased risk of preterm labor, premature rupture of membranes (PROM), chorioamnionitis, and low birth weight. Additionally, active infection at the time of delivery can result in vertical transmission of the fungus to the neonate, leading to oral thrush or diaper dermatitis.
    \item \textbf{Psychological and Psychosexual Distress}: Recurrent VVC (RVVC) has a profound negative impact on mental health. The unpredictability of episodes, chronic discomfort, and painful intercourse lead to anxiety, depression, feelings of loss of control, relationship strain, and sexual dysfunction.
    \item \textbf{Antifungal Resistance}: Inappropriate self-treatment and long-term, low-dose exposure to over-the-counter azole therapies have contributed to the selection of resistant fungal strains, particularly among non-\textit{albicans} species, making subsequent infections much harder to treat.
\end{itemize}

\section*{Prognosis and Outcomes}
The prognosis for individuals with vulvovaginal candidiasis is generally excellent. The vast majority of uncomplicated VVC cases resolve rapidly, within a few days of initiating appropriate topical or oral antifungal therapy. Patient compliance with the prescribed duration of treatment is key to achieving complete mycological clearance and preventing early relapse.

For individuals with recurrent vulvovaginal candidiasis (RVVC), the long-term prognosis is also favorable, though the clinical course can be frustratingly prolonged. Suppressive maintenance therapy with weekly fluconazole successfully controls symptoms in over 90\% of patients during the treatment period. However, recurrence remains common once the maintenance regimen is discontinued, and management often requires an ongoing partnership between the patient and their healthcare provider to identify individual triggers, optimize glycemic control in diabetic patients, and manage expectations. In postmenopausal women or those with persistent risk factors, the resolution of VVC is directly tied to the modification of the underlying predisposing factors (such as adjusting hormone replacement therapy or improving blood glucose levels).

\section*{Prevention and Self-Care}
Preventive measures and self-care strategies focus on maintaining a healthy vaginal ecosystem, minimizing local irritation, and reducing the environmental conditions that favor fungal overgrowth.

Key prevention and self-care practices include:
\begin{itemize}
    \item \textbf{Hygiene Practices}: Wash the vulvar area daily with plain water or a soap-free, pH-balanced cleanser. Avoid douching, as it strips the vagina of protective lactobacilli and disrupts the natural pH. Avoid using scented soaps, bubble baths, bath oils, feminine hygiene sprays, and scented sanitary products, which can irritate the delicate vulvar skin.
    \item \textbf{Clothing Choices}: Wear loose-fitting clothing and breathable, natural fibers, such as 100\% cotton underwear. Avoid tight jeans, synthetic underwear, tights, and pantyhose, which trap heat and moisture, creating a warm, damp environment that promotes \textit{Candida} proliferation. Change out of wet swimwear or sweaty exercise clothing as soon as possible.
    \item \textbf{Dietary and Lifestyle Modifications}: Maintain a balanced diet and limit the intake of highly refined sugars, which can elevate systemic glucose levels and promote yeast growth. For individuals with diabetes, strict glycemic control is the most critical preventive measure.
    \item \textbf{Probiotics}: While some clinical trials suggest that oral or vaginal probiotics containing \textit{Lactobacillus} species may help restore the vaginal flora, the scientific evidence supporting their efficacy in preventing or treating VVC remains mixed. They may be used as an optional supportive therapy but should not replace standard antifungal treatments.
    \item \textbf{Antibiotic Stewardship}: Use antibiotics only when prescribed by a healthcare provider for confirmed bacterial infections. If prone to VVC, discuss with a doctor the option of using a prophylactic antifungal agent (such as a single dose of oral fluconazole) at the start and completion of a course of broad-spectrum antibiotics.
    \item \textbf{Sexual Health and Lubrication}: Use water-based, glycerin-free lubricants during intercourse if vaginal dryness is present, to prevent micro-trauma to the vaginal mucosa, which can facilitate fungal invasion. Avoid lubricants containing glycerin, as it can serve as a substrate for yeast.
\end{itemize}

\section*{Sources and Further Reading}
For additional information, clinical guidelines, and patient support resources, refer to the following authoritative health organizations:
\begin{itemize}
    \item Centers for Disease Control and Prevention (CDC) -- Vulvovaginal Candidiasis Clinical Guidance: \texttt{https://www.cdc.gov/candidiasis/hcp/clinical-guidance/vulvovaginal-candidiasis.html}
    \item Mayo Clinic -- Vaginal Yeast Infection Information: \texttt{https://www.mayoclinic.org/diseases-conditions/yeast-infection/symptoms-causes/syc-20378999}
    \item The Royal Australian College of General Practitioners (RACGP) -- Vaginal Discharge Guidelines: \texttt{https://www.racgp.org.au/clinical-resources/clinical-guidelines/key-racgp-guidelines/view-all-racgp-guidelines/vaginal-discharge}
    \item British Association for Sexual Health and HIV (BASHH) -- National Guideline on the Management of Vulvovaginal Candidiasis: \texttt{https://www.bashh.org/guidelines}
    \item World Health Organization (WHO) -- Sexual and Reproductive Health and Research: \texttt{https://www.who.int/teams/sexual-and-reproductive-health-and-research}
    \item American College of Obstetricians and Gynecologists (ACOG) -- Vaginitis Clinical Information: \texttt{https://www.acog.org/womens-health/faqs/vaginitis}
\end{itemize}